% =====================================================================
%  Large Language Model Agents in Repeated Public Goods Games:
%  A Literature Review and Research Agenda
%
%  Compile:  pdflatex -> bibtex -> pdflatex -> pdflatex
%  Engine:   pdfLaTeX (MiKTeX)         Bib style: plainnat (natbib)
% =====================================================================
\documentclass[11pt,a4paper]{article}

% ---------- encoding & fonts ----------
\usepackage[T1]{fontenc}
\usepackage[utf8]{inputenc}
\usepackage{lmodern}
\usepackage{microtype}

% ---------- layout ----------
\usepackage[margin=1in]{geometry}
\usepackage{setspace}
\onehalfspacing
\usepackage{parskip}

% ---------- math ----------
\usepackage{amsmath,amssymb}

% ---------- tables & figures ----------
\usepackage{booktabs}
\usepackage{tabularx}
\usepackage{array}
\usepackage{graphicx}
\usepackage{caption}
\captionsetup{font=small,labelfont=bf}
\usepackage[edges]{forest}   % taxonomy tree (Figure 1)
\usepackage{adjustbox}       % scale the tree to fit the page
\newcolumntype{Y}{>{\raggedright\arraybackslash}X}

% ---------- lists ----------
\usepackage{enumitem}
\setlist{noitemsep,topsep=2pt}

% ---------- bibliography ----------
\usepackage[round,authoryear]{natbib}
\bibliographystyle{plainnat}

% ---------- hyperlinks ----------
\usepackage{xcolor}
% taxonomy-tree branch colours (one per major part of the review)
\definecolor{cPGG}{HTML}{C2185B}   % Background: Public Goods Game
\definecolor{cAgent}{HTML}{1565C0}  % LLMs as strategic/economic agents
\definecolor{cRep}{HTML}{2E7D32}    % LLMs in repeated games
\definecolor{cCore}{HTML}{E65100}   % LLMs in PGG & commons (core)
\definecolor{cSyn}{HTML}{6A1B9A}    % Synthesis & research agenda
\usepackage[colorlinks=true,linkcolor=blue!55!black,
            citecolor=teal!65!black,urlcolor=blue!55!black]{hyperref}

% ---------- handy macros ----------
\newcommand{\pgg}{public goods game}
\newcommand{\PGG}{Public Goods Game}
\newcommand{\llm}{LLM}
\newcommand{\eg}{e.g.,\ }
\newcommand{\ie}{i.e.,\ }

\title{\textbf{Large Language Model Agents in Repeated Public Goods Games:\\
A Literature Review and Research Agenda}}

% NOTE: replace the placeholder author/affiliation below with your own.
\author{T. C. Nguyen\\[2pt]
\small \texttt{tcnguyen2365@gmail.com}}
\date{\today}

\begin{document}
\maketitle

\begin{abstract}
\noindent
The \pgg{} (PGG) is the canonical model of the tension between individual
self-interest and collective welfare, and its repeated form is the natural
setting in which cooperation can be built, sustained, or eroded over time. As
large language models (LLMs) are increasingly deployed as autonomous, mutually
interacting agents, a fast-growing literature asks whether they cooperate,
free-ride, reciprocate, punish, and form norms the way humans do. This review
is organised around the intersection of these two threads---\emph{LLM agents in
repeated public goods games}---and is intended as an orienting map for new
research in the area. We first consolidate the classical PGG background needed
to read and benchmark agent studies: the formal model, the marginal per-capita
return, the gap between the Nash equilibrium and the social optimum, the role of
repetition, and the robust experimental and evolutionary regularities of human
play (over-contribution with decay, conditional cooperation, the power of
punishment, reputation, and institutional choice). We then survey, in depth, the
emerging body of work on LLMs as strategic and economic agents, in repeated
two-player social dilemmas, and---most directly---in public goods games and
common-pool resource dilemmas, including studies of institutional sanctioning,
emergent norms, and cooperation among populations of language models. Drawing
the threads together, we find that current LLM--PGG research is concentrated on
short horizons, a handful of models, and a narrow slice of the classical design
space, leaving the long-run dynamics that define \emph{repeated} cooperation
largely unexplored. We close with a structured analysis of open problems and a
concrete research agenda---spanning long-horizon dynamics, conditional
cooperation, the reward--punishment asymmetry, the reasoning--cooperation
trade-off, reputation, heterogeneous and mixed populations, threshold/collective-risk
variants, and evolutionary analysis of LLM-authored strategies---meant to help
identify where genuine novelty remains.

\medskip
\noindent\textbf{Keywords:} public goods game; cooperation; repeated games;
social dilemmas; large language models; multi-agent systems; free-riding;
sanctioning institutions.
\end{abstract}

% =====================================================================
\section{Introduction}
% =====================================================================

Cooperation among self-interested actors is one of the central puzzles of the
social and biological sciences \citep{nowak2006five}. Its sharpest formal
expression is the \emph{\pgg{}} (PGG): a group of individuals each privately
decide how much of an endowment to contribute to a common pool that is multiplied
and shared equally, so that the group is best off when everyone contributes while
each individual is privately best off contributing nothing
\citep{samuelson1954pure,olson1965logic,ledyard1995public}. The PGG distils the
``free-rider problem'' and the ``tragedy of the commons''
\citep{hardin1968tragedy} into a few parameters, which is why it has become the
workhorse of experimental economics, evolutionary biology, and political science
for studying when and why cooperation succeeds or collapses.

Two features make the \emph{repeated} \pgg{} especially important. First, almost
all real collective-action problems---team production, taxation, environmental
protection, open-source software, content moderation---are played out over many
interactions rather than once. Second, repetition is precisely what opens the
door to reciprocity, reputation, punishment, and norm formation: mechanisms that
can turn a one-shot dilemma with a bleak toward defection into a setting where
cooperation is an equilibrium \citep{vansegbroeck2012emergence,fehr2000cooperation}.
Decades of human experiments and evolutionary models have mapped this repeated
landscape in detail (Section~\ref{sec:background}).

A new class of decision-makers has now entered the picture. Large language models
(LLMs) are increasingly deployed not as passive tools but as \emph{autonomous
agents} that plan, negotiate, and act in multi-agent systems
\citep{park2023generative,guo2024multiagent}. When many such agents interact,
their collective behaviour becomes a first-order question for both AI safety and
the social sciences: do LLM agents cooperate for the common good, or do they
free-ride; do they reciprocate and forgive; do they enforce norms; and does this
behaviour match, exceed, or fall short of human cooperation? Because the PGG is
the standard instrument for asking exactly these questions about humans, it is a
natural and increasingly popular testbed for LLM agents as well
\citep{guzman2025corrupted,sreedhar2025simulating,piatti2024govsim}.

\paragraph{Purpose and scope.}
This article reviews the literature at the intersection of \emph{LLM agents} and
\emph{repeated public goods games}, with two goals. The first is descriptive: to
give a reader who is new to the area a coherent map of what has already been
done---both the classical PGG results that any agent study must be read against,
and the rapidly expanding set of LLM experiments in social dilemmas. The second
is generative: to synthesise these findings into an explicit list of open
problems and a research agenda, so that genuine novelty can be located rather
than rediscovered. The review is therefore deliberately weighted toward the LLM
literature (Sections~\ref{sec:agents}--\ref{sec:agenda}), while keeping a compact
but self-contained treatment of PGG theory and human evidence
(Section~\ref{sec:background}) as the foundation and the benchmark.

\paragraph{Method.}
We draw on (i) the foundational and survey literature of public goods experiments
and evolutionary game theory, and (ii) the 2022--2025 wave of papers that place
LLMs in games, social dilemmas, and multi-agent simulations. For the latter we
prioritise work that involves \emph{repeated} interaction and/or the \emph{public
goods / commons} structure, while also covering closely related two-player and
norm-formation studies that inform method and interpretation.

\paragraph{Roadmap.}
Section~\ref{sec:background} establishes the PGG background: the formal model, the
role of repetition, the main variants, and the robust regularities of human play.
Section~\ref{sec:agents} reviews LLMs as strategic and economic agents.
Section~\ref{sec:repeated} surveys LLMs in repeated two-player games and social
dilemmas. Section~\ref{sec:llmpgg}---the core of the review---examines LLMs in
public goods and common-pool resource games, including sanctioning and emergent
norms. Section~\ref{sec:synthesis} consolidates what is known and exposes the
gaps. Section~\ref{sec:agenda} proposes a concrete research agenda for LLMs in
repeated PGGs. Section~\ref{sec:conclusion} concludes.
Figure~\ref{fig:taxonomy} gives a visual taxonomy of the reviewed literature,
organised to mirror this structure.

% ---------------------------------------------------------------
% Figure 1: taxonomy tree
% ---------------------------------------------------------------
\begin{figure}[p]
\centering
\begin{adjustbox}{max width=\textwidth,max totalheight=0.85\textheight,center}
\begin{forest}
  forked edges,
  for tree={
    grow'=east,
    draw=black!50,
    rounded corners=2pt,
    line width=0.5pt,
    node options={align=left},
    inner xsep=4pt,
    inner ysep=3pt,
    anchor=west,
    edge={draw=black!45, line width=0.5pt},
    l sep=8pt,
    s sep=3pt,
    font=\footnotesize,
    text width=3.0cm,                                   % default: every node wraps
    if n children=0{text width=5.4cm, font=\scriptsize}{}, % leaves: wider, smaller
  },
  where level=0{text width=2.6cm, font=\small\bfseries, fill=black!8}{},
  [{LLM Agents in Repeated Public Goods Games}
    %------------------ Branch 1 : PGG background ------------------
    [{Background:\\Public Goods Game (\S\ref{sec:background})},
        fill=cPGG!32, draw=cPGG!60!black,
        for descendants={fill=cPGG!10, draw=cPGG!55!black}
      [{Model \&\\theory}
        [{Payoff, MPCR, Nash vs.\ optimum \citep{samuelson1954pure,olson1965logic,ledyard1995public}}]
        [{Repetition \& folk theorem \citep{camerer2003behavioral}}]
      ]
      [{Variants}
        [{Punishment / reward \citep{fehr2000cooperation,andreoni2003carrot}}]
        [{Institutional choice \citep{gurerk2006competitive}}]
        [{Threshold / collective-risk \citep{milinski2008collective,santos2011risk}}]
        [{Common-pool resource \citep{ostrom1990governing}}]
      ]
      [{Human\\regularities}
        [{Over-contribution \& decay \citep{isaac1988groupsize,andreoni1988free}}]
        [{Conditional cooperation \citep{fischbacher2001conditional,fischbacher2010social}}]
        [{Punishment \& reputation \citep{fehr2002altruistic,milinski2002reputation}}]
      ]
      [{Evolutionary\\dynamics}
        [{Replicator, networks, diversity \citep{hofbauer1998evolutionary,santos2008social}}]
        [{Repeated fairness, risk \citep{vansegbroeck2012emergence,santos2011risk}}]
      ]
    ]
    %------------------ Branch 2 : strategic/economic agents ------------------
    [{LLMs as Strategic \&\\Economic Agents (\S\ref{sec:agents})},
        fill=cAgent!30, draw=cAgent!60!black,
        for descendants={fill=cAgent!10, draw=cAgent!55!black}
      [{Simulated economic agents \citep{horton2023homosilicus,mei2024turing,brookins2024playing}}]
      [{Strategic reasoning \& limits \citep{fan2024rational,lore2024strategic,duan2024gtbench,hua2024gametheoretic}}]
      [{Generative-agent architectures \citep{park2023generative,li2024econagent}}]
    ]
    %------------------ Branch 3 : repeated games ------------------
    [{LLMs in Repeated Games \&\\Social Dilemmas (\S\ref{sec:repeated})},
        fill=cRep!30, draw=cRep!60!black,
        for descendants={fill=cRep!10, draw=cRep!55!black}
      [{Repeated 2-player / IPD \citep{akata2025playing,fontana2025nicer,willis2025cooperate,tonini2025superadditive}}]
      [{Trust \& communication \citep{xie2024trust,xu2023werewolf}}]
    ]
    %------------------ Branch 4 : PGG & commons (core) ------------------
    [{LLMs in Public Goods\\Games \& the Commons (\S\ref{sec:llmpgg})},
        fill=cCore!38, draw=cCore!65!black,
        for descendants={fill=cCore!13, draw=cCore!60!black}
      [{Replicating \& benchmarking human PGG \citep{bao2023cognitive,sreedhar2025simulating,xie2024chatbots}}]
      [{Reasoning, free-riding, sanctioning \citep{guzman2025corrupted,li2025spontaneous,long2025mirror,madmoun2025communication}}]
      [{Commons \& sustainability \citep{piatti2024govsim,mosquera2024meltingpot,nguyen2025navigating,gupta2025social}}]
      [{Threshold / collective-risk \citep{slumbers2025port,willis2026hundreds}}]
      [{Reputation, norms \& populations \citep{leng2023llmsocial,horiguchi2024evolution,vallinder2024cultural,ashery2025emergent,cross2025validating,zhu2026gossip,liu2026memory}}]
    ]
    %------------------ Branch 5 : synthesis & agenda ------------------
    [{Synthesis \&\\Research Agenda (\S\ref{sec:synthesis}--\ref{sec:agenda})},
        fill=cSyn!30, draw=cSyn!60!black,
        for descendants={fill=cSyn!10, draw=cSyn!55!black}
      [{Established: cooperate partly like humans; model/prompt-dependent; reasoning $\downarrow$ cooperation}]
      [{Gaps: short horizons; conditional cooperation; reward vs.\ punishment; reputation}]
      [{Agenda: long-horizon benchmark; threshold/risk PGG; evolutionary analysis; mixed populations}]
    ]
  ]
\end{forest}
\end{adjustbox}
\caption{A taxonomy of the literature on LLM agents in repeated public goods
games, mirroring the structure of this review. Branch colours denote the major
parts: the PGG background (Section~\ref{sec:background}), LLMs as strategic and
economic agents (Section~\ref{sec:agents}), LLMs in repeated games and social
dilemmas (Section~\ref{sec:repeated}), LLMs in public goods and commons games
(Section~\ref{sec:llmpgg}, the core), and the synthesis and research agenda
(Sections~\ref{sec:synthesis}--\ref{sec:agenda}). Leaf nodes list representative
references.}
\label{fig:taxonomy}
\end{figure}

% =====================================================================
\section{Background: The Public Goods Game}
\label{sec:background}
% =====================================================================

This section provides the conceptual grounding needed to interpret and benchmark
LLM studies. Readers familiar with the PGG may skim to
Section~\ref{sec:agents}; the stylized facts in Table~\ref{tab:facts} are the
human baseline against which agent behaviour is repeatedly compared in later
sections.

% ---------------------------------------------------------------
\subsection{The linear public goods game}
% ---------------------------------------------------------------

In the standard linear PGG, $N$ players each receive an endowment $e$ and
simultaneously choose a contribution $c_i\in[0,e]$ to a common pool. The pool is
multiplied by a factor $\alpha$ (with $1<\alpha<N$) and divided equally. Player
$i$'s payoff is
\begin{equation}
\pi_i \;=\; \underbrace{(e-c_i)}_{\text{kept}} \;+\;
\frac{\alpha}{N}\sum_{j=1}^{N} c_j .
\label{eq:pgg}
\end{equation}
The ratio $\mathrm{MPCR}=\alpha/N$ is the \emph{marginal per-capita return}: the
private return to each unit contributed. Because $\alpha<N$ implies
$\mathrm{MPCR}<1$, every unit a player contributes yields them less than a unit
kept, so contributing nothing is the dominant strategy and $c_i=0$ for all $i$ is
the unique Nash equilibrium. Yet because $\alpha>1$, total welfare is maximised
when everyone contributes fully ($c_i=e$). The wedge between the individually
rational outcome (universal free-riding) and the Pareto optimum (full
contribution) is the social dilemma in its purest form
\citep{ledyard1995public,camerer2003behavioral}.

% ---------------------------------------------------------------
\subsection{Why repetition matters}
% ---------------------------------------------------------------

In a \emph{one-shot} PGG, or in a \emph{finitely} repeated PGG solved by backward
induction, the theoretical prediction is unrelenting defection. The picture
changes once interaction is genuinely repeated. In an \emph{infinitely} (or
indefinitely) repeated game, the folk theorem implies that cooperation can be
sustained as a subgame-perfect equilibrium when players are patient enough,
because future cooperation can be made contingent on present behaviour through
strategies such as conditional cooperation and trigger strategies
\citep{camerer2003behavioral}. Repetition is thus not a technical detail but the
very thing that makes reciprocity, reputation, punishment, and norm formation
possible. It is also where human behaviour departs most strikingly from naive
theory (Section~\ref{subsec:facts}), and---as this review emphasises---where the
behaviour of LLM agents is least understood.

% ---------------------------------------------------------------
\subsection{Variants relevant to agent studies}
% ---------------------------------------------------------------

Several PGG variants recur in the agent literature and are worth naming:
\begin{itemize}
  \item \textbf{PGG with punishment / reward.} After contributions are revealed,
  players may pay a cost to reduce (punish) or increase (reward) another's payoff.
  Costly peer punishment dramatically raises cooperation in human experiments
  \citep{fehr2000cooperation,fehr2002altruistic}, and rewards and punishments have
  distinct effects \citep{andreoni2003carrot}.
  \item \textbf{Institutional / sanctioning choice.} Players choose between a
  sanction-free institution and one that permits punishment; over time people
  migrate to the sanctioning institution, which out-competes its rival
  \citep{gurerk2006competitive}. This design is directly adapted in recent
  LLM work \citep{guzman2025corrupted}.
  \item \textbf{Threshold / collective-risk PGG.} A collective benefit (or the
  avoidance of a loss) is realised only if total contributions reach a threshold;
  otherwise players risk losing their endowment. This captures climate-change-type
  dilemmas \citep{milinski2008collective,santos2011risk} and is discussed in
  Section~\ref{subsec:evo}.
  \item \textbf{Common-pool resource (CPR) / commons.} The ``inverse'' of the
  contribution PGG: players \emph{extract} from a shared, regenerating resource,
  and over-extraction causes collapse \citep{ostrom1990governing}. LLM versions
  of this dilemma are prominent \citep{piatti2024govsim}.
  \item \textbf{Structured populations.} Contributions interact through networks
  or are subject to evolutionary updating; heterogeneity and network structure
  strongly shape cooperation \citep{santos2008social,szabo2007evolutionary}.
\end{itemize}

% ---------------------------------------------------------------
\subsection{Robust regularities of human play}
\label{subsec:facts}
% ---------------------------------------------------------------

Five decades of laboratory work have produced a set of remarkably robust
``stylized facts'' (Table~\ref{tab:facts}). People do \emph{not} free-ride
completely: initial contributions in linear PGGs typically average around
40--60\% of the endowment \citep{ledyard1995public,zelmer2003linear}. But under
repetition without any cooperation-supporting mechanism, contributions
\emph{decay} toward (though rarely reaching) the free-riding prediction
\citep{isaac1988groupsize,andreoni1988free}. The decay is best explained not by
pure confusion---though confusion contributes \citep{andreoni1995kindness,
houser2002revisiting}---but by the heterogeneity of types and \emph{imperfect
conditional cooperation}: roughly half of subjects raise their contribution when
others do, but by slightly less than matching, so beliefs and contributions
ratchet downward over time \citep{fischbacher2001conditional,fischbacher2010social,
croson2007theories}. Framing, group size and the MPCR, communication, reputation,
and---above all---the opportunity to punish free-riders are the main levers that
move cooperation \citep{andreoni1995warmglow,isaac1994groupsize,
fehr2000cooperation,milinski2002reputation,chaudhuri2011sustaining}. These
regularities are the natural benchmark for any claim that LLM agents do or do not
``behave like humans''.

\begin{table}[t]
\centering
\small
\caption{Robust regularities of human behaviour in (repeated) public goods
games. These constitute the empirical baseline against which LLM agents are
compared in Sections~\ref{sec:llmpgg}--\ref{sec:synthesis}.}
\label{tab:facts}
\begin{tabularx}{\textwidth}{@{}p{0.30\textwidth} Y@{}}
\toprule
\textbf{Regularity} & \textbf{Summary and key references} \\
\midrule
Partial cooperation, not full free-riding &
Initial contributions average $\sim$40--60\% of endowment, far above the Nash
prediction of zero \citep{ledyard1995public,zelmer2003linear}. \\
Decay under repetition &
Without supporting mechanisms, contributions decline toward zero over rounds
\citep{isaac1988groupsize,andreoni1988free}. \\
Conditional cooperation &
$\sim$50\% of people are conditional cooperators who match others imperfectly;
$\sim$30\% are free-riders \citep{fischbacher2001conditional,fischbacher2010social}. \\
Kindness \emph{and} confusion &
Cooperation reflects genuine other-regard plus decision error in roughly equal
measure \citep{andreoni1995kindness,houser2002revisiting}. \\
Framing effects &
Positive (``warm-glow'') vs.\ negative framing changes cooperation for identical
payoffs \citep{andreoni1995warmglow}. \\
MPCR and group size &
Higher MPCR raises contributions; group-size effects are largely an MPCR artefact
\citep{isaac1988groupsize,isaac1994groupsize}. \\
Punishment sustains cooperation &
Costly peer punishment sharply raises and stabilises contributions, even when
individually costly \citep{fehr2000cooperation,fehr2002altruistic}. \\
Reputation and institutions &
Indirect reciprocity and a choice of sanctioning institutions both rescue
cooperation \citep{milinski2002reputation,gurerk2006competitive,rockenbach2006efficient}. \\
\bottomrule
\end{tabularx}
\end{table}

% ---------------------------------------------------------------
\subsection{Repeated-group and evolutionary dynamics}
\label{subsec:evo}
% ---------------------------------------------------------------

Beyond the laboratory, evolutionary game theory studies how cooperation fares
when strategies are updated over time within a population, typically through the
replicator dynamics $\dot{x}=x(1-x)(f_C-f_D)$, where $x$ is the fraction of
cooperators and $f_C,f_D$ their average payoffs
\citep{hofbauer1998evolutionary,sigmund2010calculus}. In well-mixed populations
cooperation tends to vanish, but a battery of mechanisms can reverse this:
network reciprocity and social diversity \citep{santos2008social,szabo2007evolutionary,
perc2017statistical}, optional participation \citep{hauert2002volunteering},
the emergence of costly punishment \citep{hauert2007freedom,boyd2003evolution},
institutions learned through social imitation \citep{sigmund2010social}, and
multilevel selection \citep{traulsen2006evolution}; \citet{nowak2006five}
organises these into five canonical rules.

Two strands are especially relevant to \emph{repeated} PGGs and to the LLM agenda.
First, \citet{vansegbroeck2012emergence} model the \emph{repeated} $N$-person
game directly: agents are ``reciprocators'' who contribute only if at least $M$
others contributed in the previous round, and evolution selects for an
intermediate, moderate notion of fairness $M^\ast$, producing cyclic alternations
between cooperation and defection rather than a static equilibrium. Second, the
\emph{collective-risk} dilemma combines repetition with a threshold and the risk
of catastrophic loss: in \citeauthor{milinski2008collective}'s
\citeyearpar{milinski2008collective} climate experiment, groups reached the target
only when the perceived risk of loss was high, and \citet{santos2011risk} show
analytically that a high risk of collective failure---together with small group
size and stringent thresholds---is what turns cooperation from a losing strategy
into a viable one. These results give concrete, quantitative hypotheses (the role
of risk, thresholds, group size, and reciprocation rules) that LLM studies have
barely begun to test.

% =====================================================================
\section{LLMs as Strategic and Economic Agents}
\label{sec:agents}
% =====================================================================

Before turning to multi-agent public goods settings, it is useful to summarise
what is known about individual LLM behaviour in economic and strategic tasks,
because these capabilities and biases propagate into the PGG.

\paragraph{LLMs as simulated economic agents.}
\citet{horton2023homosilicus} argued that an LLM endowed with a persona can serve
as ``\emph{homo silicus}''---a simulated economic agent whose responses reproduce
classic behavioural results, suggesting LLMs can act as in-silico subjects for
economics. \citet{mei2024turing} formalised this with a behavioural Turing test
across a battery of personality measures and games (dictator, trust, public goods,
prisoner's dilemma), finding that GPT-4 generally behaves within the human
distribution and is often \emph{more} cooperative and altruistic than the median
person. \citet{brookins2024playing} similarly found GPT to display stronger
fairness and higher cooperation than typical human benchmarks in dictator and
prisoner's dilemma games. The recurring theme is partial human-likeness with a
prosocial tilt---and substantial sensitivity to how the situation is framed.

\paragraph{Strategic reasoning: capability and limits.}
A parallel line evaluates whether LLMs reason strategically in the game-theoretic
sense. Results are mixed. \citet{fan2024rational} report that even GPT-4 deviates
systematically from rational benchmarks in dictator, rock--paper--scissors, and
ring-network games. \citet{lore2024strategic} disentangle the effects of formal
game structure from contextual framing across the prisoner's dilemma, stag hunt,
and snowdrift, showing that LLMs are swayed by surface framing more than humans.
\citet{duan2024gtbench} introduce \textsc{GTBench}, a suite of game-theoretic
tasks, and find LLMs strong on probabilistic, incomplete-information games but
weak on deterministic, complete-information ones. \citet{hua2024gametheoretic}
show that deviations from equilibrium grow with game complexity and propose
explicit ``game-theoretic workflows'' to improve equilibrium computation, while
\citet{gandhi2023strategic} demonstrate that structured chain-of-thought prompting
(tracking states, values, and beliefs) markedly improves strategic play. The
implication for PGGs is twofold: LLM cooperation is real but is mediated by
prompting and framing, and ``better reasoning'' does not automatically translate
into better strategic---or more cooperative---behaviour.

\paragraph{Generative-agent architectures.}
Methodologically, much multi-agent PGG work builds on the generative-agent
paradigm of \citet{park2023generative}, in which agents are equipped with memory,
reflection, and planning modules and exhibit believable emergent social behaviour
in a sandbox. \citet{li2024econagent} extend this to macroeconomic simulation
with perception/memory/reflection agents. These architectures supply the
machinery---persistent memory, identities, and natural-language interaction---that
distinguishes a multi-agent LLM PGG from a sequence of isolated prompts, and that
will be central to studying \emph{repeated} cooperation.

% =====================================================================
\section{LLMs in Repeated Games and Social Dilemmas}
\label{sec:repeated}
% =====================================================================

The closest precedent to repeated PGGs is the study of LLMs in repeated
two-player social dilemmas, where the dynamics of reciprocity and forgiveness can
be observed cleanly.

\paragraph{Repeated two-player games.}
\citet{akata2025playing} systematically played families of repeated $2\times2$
games with LLMs and found a characteristic profile: LLMs coordinate well but are
\emph{unforgiving} in the iterated prisoner's dilemma, defecting persistently
after a single defection by the opponent---behaviour that is individually
defensible but corrosive to long-run cooperation. \citet{fontana2025nicer}, by
contrast, report that several models are ``nicer than humans'' in the iterated
prisoner's dilemma, though newer models can be more exploitative, underscoring
that conclusions are model-dependent. \citet{phelps2023investigating} use
experimental-economics prompts to probe emergent goal-like dispositions and show
that persona and framing strongly steer cooperation. \citet{willis2025cooperate}
take an evolutionary view: they prompt LLMs to \emph{author} full iterated-PD
strategies and then analyse those strategies with evolutionary game theory,
revealing systematic, model-specific biases toward aggression or cooperation.
\citet{tonini2025superadditive} show that pairing within-team repetition with
between-team competition yields ``super-additive'' cooperation, raising even
one-shot cooperation rates---an LLM echo of multilevel selection.

\paragraph{Trust, communication, and other multi-agent games.}
\citet{xie2024trust} find, using behavioural trust games, that LLM agents exhibit
trust behaviour aligning closely with humans for GPT-4. Communication-rich games
such as Werewolf reveal emergent strategic behaviours (trust, bluffing,
leadership) from frozen models \citep{xu2023werewolf}. Together these studies
establish that LLM agents can reciprocate, trust, and communicate strategically
over repeated interactions---but also that they carry strong, idiosyncratic, and
prompt-sensitive dispositions that must be controlled for when moving to the
$N$-player public goods setting.

% =====================================================================
\section{LLMs in Public Goods Games and the Commons}
\label{sec:llmpgg}
% =====================================================================

We now reach the core of the review: studies that place LLM agents directly in
public goods or common-pool resource dilemmas. Table~\ref{tab:llm} summarises the
key studies along the design dimensions that matter most for repeated cooperation.

% ---------------------------------------------------------------
\subsection{Replicating human PGG behaviour}
% ---------------------------------------------------------------

The earliest direct placements of an LLM in a PGG simply elicited contributions:
\citet{bao2023cognitive} had GPT play a voluntary-contribution game and linked its
contribution level to cognitive uncertainty. Behavioural-economics benchmark
batteries now routinely include the PGG among a suite of canonical games
\citep{mei2024turing,xie2024chatbots,leng2023llmsocial}, generally finding that
capable LLMs contribute at or above human levels, but with pronounced variation
across model families.

\citet{sreedhar2025simulating} ask whether multi-agent LLM systems can reproduce
human prosocial behaviour in the PGG. Using a generative-agent architecture with
three or four player agents, they replicate three classic human treatments---name
framing (``Teamwork'' vs.\ ``Taxation''), transparency of contributions, and
varying endowments---and find that LLM agents reproduce the \emph{direction} of
each human effect (\eg priming toward teamwork raises contributions; transparency
raises contributions; large endowments lower the contributed fraction), although
they tend to \emph{overestimate the magnitude}. They further show that effects can
transfer from non-PGG experiments and that, with added communication channels,
agents display ``unbounded'' real-world behaviours such as collaborating and even
cheating. The study is important evidence that LLMs encode a usable model of human
cooperative tendencies, and that this model is steerable by the same levers
(framing, transparency, endowment) that move humans. Its limits---mostly one-shot
or short (five-round) games, a single primary model, no sanctioning---also mark
where repeated-PGG work must go.

% ---------------------------------------------------------------
\subsection{Reasoning, free-riding, and institutional sanctioning}
% ---------------------------------------------------------------

The work closest to a genuine \emph{repeated PGG with institutions} is
\citet{guzman2025corrupted}. Adapting the institutional-choice design of
\citet{gurerk2006competitive}, they run a 15-round PGG with seven LLM agents who
each round choose between a \emph{sanctioning institution} (SI), in which members
may pay to reward (\(+1\) at cost \(1\)) or punish (\(-3\) at cost \(1\)) one
another, and a \emph{sanction-free institution} (SFI); contributions are
multiplied by 1.6 and shared within the chosen institution. Their findings are
striking and directly relevant:
\begin{itemize}
  \item \textbf{Four behavioural archetypes} emerge across models---\emph{increasingly
  cooperative}, \emph{increasingly defecting}, \emph{no change}, and
  \emph{unstable (oscillating)}---closely mirroring the heterogeneity seen in
  human and evolutionary studies.
  \item \textbf{Reasoning models cooperate \emph{less}.} Traditional models
  (\eg GPT-4o, DeepSeek-V3, Llama-3.3-70B) sustain high cooperation comparable to
  humans, whereas reasoning-optimised models (the o1/o3 family) frequently
  collapse to free-riding---``corrupted by reasoning''. This is a pointed
  counterexample to the assumption that stronger reasoning yields better
  collective outcomes.
  \item \textbf{Reward, not punishment.} Although capable LLMs replicate human
  cooperative \emph{outcomes}, they do so through a different \emph{mechanism}:
  humans rely heavily on punishment ($\sim$1.66 punish/reward ratio), whereas all
  LLMs strongly prefer rewards, raising doubts about whether they grasp deterrence
  and the stability of the cooperation they achieve.
\end{itemize}
This study is the clearest demonstration to date that the \emph{repeated},
institution-rich PGG is a revealing probe of LLM cooperation, and that
architecture (reasoning vs.\ not) is a first-order determinant of outcomes.

The headline result---that reasoning optimisation can \emph{erode}
cooperation---has since been corroborated and extended. \citet{li2025spontaneous}
find, across the PGG and five other cooperation games, that eliciting
chain-of-thought ``calculated'' reasoning systematically lowers giving and norm
enforcement relative to more ``spontaneous'' responses, and
\citet{backmann2025moralsim} (MoralSim) show that no frontier model stays
cooperative once payoffs conflict with stated moral norms in repeated prisoner's
dilemma and PGG play. \citet{long2025mirror} add a twist: in an iterated PGG,
telling a model that its co-player is ``another AI'' versus ``itself''
measurably shifts contributions. A complementary line uses the repeated PGG
constructively: \citet{madmoun2025communication} show that cheap-talk
communication restores cooperation in an iterated PGG with punishment (while a
rationality-maximising training curriculum suppresses it),
\citet{warnakulasuriya2025evolution} replicate punishment-driven norm emergence in
an $N$-player sanctioning dilemma, and \citet{liang2025everyone} redesign a
sequential PGG as a mechanism that makes contribution the unique equilibrium for
multi-LLM ensembles. The consistent thread is that, as for humans, LLM cooperation
in the repeated PGG is fragile to ``rational'' framing and strongly helped by
communication and sanctions.

% ---------------------------------------------------------------
\subsection{Common-pool resources and sustainability}
% ---------------------------------------------------------------

\citet{piatti2024govsim} introduce \textsc{GovSim}, a common-pool resource
simulation (fishery/pasture/pollution variants) in which LLM agents repeatedly
harvest a regenerating resource. Most models over-harvest and drive the resource
to collapse; only the strongest models, and only when allowed to \emph{communicate}
and to engage in ``universalisation'' moral reasoning (``what if everyone did
this?''), achieve sustainable cooperation. \textsc{GovSim} is the commons (or
``inverse'') counterpart to the contribution PGG and demonstrates two themes that
recur throughout this review: communication is a powerful lever for LLM
cooperation, and sustainability over long horizons is fragile and
model-dependent.

The commons side has since become comparatively well populated.
\citet{mosquera2024meltingpot} were among the first to place LLM-augmented agents
in DeepMind's \textsc{Melting Pot} ``Commons Harvest'' substrate, and
\citet{chaconchamorro2025resilience} use the same family of environments to
compare the \emph{cooperative resilience} of humans and LLMs under stochastic
disruptions. Building on \textsc{GovSim}, \citet{nguyen2025navigating} equip
agents with an explicit ``consideration of future consequences'' to sustain shared
resources, while \citet{gupta2025social} embed Ostrom-style social learning and
norm-based punishment so that conservation norms emerge endogenously without reward
shaping. Across these studies the message echoes \citet{piatti2024govsim}:
collapse is the default outcome, and communication, foresight, and norms are what
avert it.

% ---------------------------------------------------------------
\subsection{Threshold and collective-risk dilemmas}
% ---------------------------------------------------------------

The threshold / collective-risk variant---contribute enough to avert a
catastrophic loss, or risk losing everything (Section~\ref{subsec:evo})---is the
LLM niche most directly connected to the climate-style dilemmas of
\citet{milinski2008collective} and \citet{santos2011risk}, and it remains almost
empty. The closest work is \citet{slumbers2025port}, who have LLM agents play
\emph{Port of Mars}, a multi-round collective-risk social dilemma with a
catastrophe threshold; however, it is framed as a test of whether LLMs can
substitute for human experimental subjects rather than as a study of
threshold-cooperation dynamics, and it uses a single model.
\citet{willis2026hundreds} include a collective-risk dilemma among three
population-scale games, but in a paradigm where models \emph{author} strategies
that are then executed automatically, rather than deliberating round by round. A
dedicated, multi-model study of LLM agents \emph{playing} a repeated
collective-risk PGG---probing risk sensitivity, threshold behaviour, group size,
and framing---does not yet appear to exist, a gap we return to in
Section~\ref{sec:agenda}.

% ---------------------------------------------------------------
\subsection{Reputation, norms, and populations of agents}
% ---------------------------------------------------------------

A further cluster studies cooperation and norms among \emph{populations} of LLM
agents. \citet{leng2023llmsocial} run canonical behavioural-economics games
(including the PGG) and show that LLM reasoning invoking fairness or altruism
predicts prosocial choices. \citet{horiguchi2024evolution} find that agents
conversing in natural language can spontaneously develop social norms and even
\emph{metanorms} (punishing those who fail to punish), echoing the second-order
free-rider problem central to human cooperation theory.
\citet{vallinder2024cultural} study indirect reciprocity across
\emph{generations} of agents in an iterated Donor Game and find that cooperation
emerges for some model families (\eg Claude) but not others, and is shaped by the
availability of punishment and by cultural transmission. \citet{ashery2025emergent}
show that decentralised LLM populations spontaneously converge on shared
conventions, with tipping points driven by committed minorities. These studies
collectively show that the ingredients of repeated-PGG cooperation---reciprocity,
reputation, punishment, norms---can all emerge among LLM agents, but unevenly and
with strong model and design dependence.

Two further threads sharpen this picture. On \emph{reputation},
\citet{cross2025validating} show that generative LLM agents reproduce the
empirical boost that gossip and ostracism give to PGG contributions, and
\citet{zhu2026gossip} build a decentralised gossip network in which self-interested
agents form reputations and ostracise defectors---directly importing indirect
reciprocity into the LLM PGG. On \emph{scale and horizon},
\citet{willis2026hundreds} push social-dilemma simulations to hundreds of agents
(finding that newer models trend toward \emph{worse} collective equilibria),
\citet{liu2026memory} run social dilemmas including PGG-type games for 500 rounds
and report a ``memory curse'' in which longer recall \emph{erodes} cooperation, and
\citet{tewolde2026coopeval} benchmark cooperation-sustaining mechanisms
(repetition, reputation, mediation, contracts) across prisoner's dilemma and
public-goods settings.

\begin{table}[t]
\centering
\footnotesize
\caption{Representative studies of LLM agents in public goods, commons, and
closely related social dilemmas, along the design dimensions most relevant to
\emph{repeated} cooperation. ``Rounds'' indicates the degree of repetition;
``Sanction'' indicates whether punishment/reward or institutional choice is
available; ``Comm.'' indicates free-form natural-language communication.}
\label{tab:llm}
\begin{tabularx}{\textwidth}{@{}p{0.17\textwidth} p{0.15\textwidth} c c c Y@{}}
\toprule
\textbf{Study} & \textbf{Game} & \textbf{Rounds} & \textbf{Sanc.} &
\textbf{Comm.} & \textbf{Key finding} \\
\midrule
\citet{sreedhar2025simulating} & PGG (contribution) & 1--5 & No & Partly &
LLMs reproduce the direction of human framing/transparency/endowment effects, but
overshoot the magnitude. \\
\citet{guzman2025corrupted} & PGG + institutional choice & 15 & Yes & No &
Reasoning models free-ride; traditional models cooperate; all LLMs prefer reward
over punishment. \\
\citet{li2025spontaneous} & PGG (+5 coop.\ games) & Many & Yes & No &
Chain-of-thought ``calculated'' reasoning lowers giving and norm enforcement. \\
\citet{piatti2024govsim} & Common-pool resource & Many & No & Yes &
Most models over-harvest and collapse; communication + universalisation enable
sustainability. \\
\citet{slumbers2025port} & Collective-risk (Port of Mars) & 9 & No & Yes &
LLMs play a catastrophe-threshold dilemma; tested as human-subject surrogates. \\
\citet{nguyen2025navigating} & Common-pool resource & Many & No & Yes &
``Consideration of future consequences'' sustains the shared resource. \\
\citet{leng2023llmsocial} & PGG, dictator, PD & 1 & No & No &
Reasoning that cites fairness/altruism predicts prosocial LLM actions. \\
\citet{horiguchi2024evolution} & Norms game & Many & Yes & Yes &
Agents develop norms and metanorms (punishing non-punishers) in natural language. \\
\citet{vallinder2024cultural} & Donor game (indirect recip.) & Many (gen.) & Yes &
No & Cooperation emerges for some model families; shaped by punishment and cultural
transmission. \\
\citet{akata2025playing} & Repeated $2\times2$ (incl.\ IPD) & Many & No & No &
LLMs coordinate well but are unforgiving in the IPD. \\
\citet{willis2025cooperate} & Iterated PD (evolutionary) & Many & No & No &
LLM-authored strategies show model-specific aggression/cooperation biases. \\
\citet{tonini2025superadditive} & IPD + group competition & Many & No & No &
Within-team repetition $+$ between-team competition yields super-additive
cooperation. \\
\citet{cross2025validating} & PGG + gossip/ostracism & Many & Yes & Yes &
Generative agents reproduce reputation and punishment effects on contributions. \\
\citet{liu2026memory} & Social dilemmas (incl.\ PGG) & 500 & No & No &
``Memory curse'': longer recall erodes cooperative intent over long horizons. \\
\bottomrule
\end{tabularx}
\end{table}

% =====================================================================
\section{Synthesis: What Is Established, and What Is Missing}
\label{sec:synthesis}
% =====================================================================

% ---------------------------------------------------------------
\subsection{What the literature has established}
% ---------------------------------------------------------------

Reading the agent studies against the human baseline of
Section~\ref{sec:background}, several conclusions are now reasonably well
supported:
\begin{enumerate}
  \item \textbf{LLMs can cooperate, and partly like humans.} Across PGGs and
  related games, capable LLMs contribute substantially and reproduce the
  \emph{direction} of major human effects (framing, transparency, endowment,
  benefits of communication) \citep{sreedhar2025simulating,mei2024turing,
  piatti2024govsim}.
  \item \textbf{Behaviour is strongly model- and prompt-dependent.} Cooperation
  varies sharply across model families and is sensitive to framing and persona
  \citep{fontana2025nicer,lore2024strategic,vallinder2024cultural}; conclusions
  rarely generalise across all models.
  \item \textbf{Reasoning optimisation can reduce cooperation.} The
  reasoning-tuned o1/o3 family free-rides more than standard models in the
  repeated institutional PGG \citep{guzman2025corrupted}, and ``calculated''
  chain-of-thought lowers giving across the PGG and related games
  \citep{li2025spontaneous,backmann2025moralsim}---complicating the assumption that
  capability implies prosociality.
  \item \textbf{Mechanism, not just outcome, differs from humans.} LLMs lean on
  rewards rather than punishment \citep{guzman2025corrupted} and can be
  unforgiving after a single defection \citep{akata2025playing}---outcomes can
  match humans while the underlying strategy does not.
  \item \textbf{The ingredients of repeated cooperation can emerge.}
  Reciprocity, reputation, punishment, norms, and even metanorms appear among LLM
  agents \citep{horiguchi2024evolution,vallinder2024cultural,ashery2025emergent},
  but unevenly.
\end{enumerate}

% ---------------------------------------------------------------
\subsection{Methodological gaps}
% ---------------------------------------------------------------

The field is young, and its methods are correspondingly narrow:
\begin{itemize}
  \item \textbf{Short horizons.} Most ``repeated'' LLM PGGs run only a handful to
  a few tens of rounds (Table~\ref{tab:llm}); only isolated recent work reaches
  hundreds of rounds \citep{liu2026memory}. The long-run phenomena that
  \emph{define} repeated cooperation---steady-state contribution decay, endgame
  effects, the slow build-up of reputation, cyclic cooperation/defection
  \citep{vansegbroeck2012emergence}---remain largely untested.
  \item \textbf{Few models, weak statistics.} Studies often use one or two
  proprietary models, sometimes with single runs per condition, limiting
  statistical power and reproducibility.
  \item \textbf{Narrow design slice.} Communication, persistent identity,
  reputation, threshold/risk structures, and large groups are each present in only
  a subset of studies, and rarely in combination.
  \item \textbf{No standard benchmark.} There is no agreed-upon repeated-PGG
  protocol, set of metrics, or human-baseline comparison, making cross-study
  synthesis difficult.
\end{itemize}

% ---------------------------------------------------------------
\subsection{Conceptual gaps mapped to the classical literature}
% ---------------------------------------------------------------

More importantly, mapping agent studies onto the human/evolutionary PGG
literature exposes specific, high-value questions that remain open:
\begin{itemize}
  \item \textbf{Conditional cooperation.} Are LLMs conditional cooperators in the
  precise sense of \citet{fischbacher2001conditional}---and do they exhibit the
  same imperfect, self-serving matching that drives human decay? A strategy-method
  analog for LLMs is largely missing.
  \item \textbf{Contribution decay.} Do LLM contributions decay over long
  horizons as human contributions do \citep{isaac1988groupsize,fischbacher2010social},
  and is any decay driven by belief updating or by context-window/memory
  effects---as the ``memory curse'' of \citet{liu2026memory} begins to suggest?
  \item \textbf{Reward--punishment asymmetry.} Why do LLMs prefer reward to
  punishment \citep{guzman2025corrupted}, is it robust across designs, and does it
  create a second-order free-rider problem in enforcement?
  \item \textbf{The reasoning--cooperation trade-off.} What mechanism makes
  reasoning models free-ride---chain-of-thought toward ``rational'' defection,
  RLHF objectives, or framing? Can alignment or prompting restore cooperation
  without harming capability?
  \item \textbf{Reputation and indirect reciprocity.} Human cooperation is
  rescued by reputation \citep{milinski2002reputation,rockenbach2006efficient},
  yet leading LLM PGGs \emph{anonymise} agents each round
  \citep{guzman2025corrupted}. First steps re-introduce gossip and reputation
  \citep{cross2025validating,zhu2026gossip}, but persistent identities and indirect
  reciprocity in the institutional PGG remain largely open.
  \item \textbf{Heterogeneous and mixed populations.} Almost all studies use
  homogeneous model populations. Mixed-model societies, and LLM--human mixed
  groups, are barely explored despite their practical importance.
  \item \textbf{Threshold/collective-risk PGGs.} The risk-driven cooperation of
  \citet{milinski2008collective} and \citet{santos2011risk} has only just begun to
  be probed with LLMs---through a single-model collective-risk simulation
  \citep{slumbers2025port} and a strategy-authoring population study
  \citep{willis2026hundreds}. Whether LLM agents \emph{deliberating} round by round
  are risk-sensitive in the way that rescues human cooperation, and how this varies
  across models and framings, remains essentially unknown.
  \item \textbf{Evolutionary/population analysis.} \citet{willis2025cooperate}
  analyse LLM-authored \emph{IPD} strategies with evolutionary game theory, and
  \citet{willis2026hundreds} extend the strategy-authoring paradigm to PGG,
  collective-risk, and commons games at population scale; a complementary analysis
  of LLM agents \emph{playing} the $N$-person PGG turn by turn (their ESS,
  replicator dynamics, invasion by free-riders) is still open.
\end{itemize}

% =====================================================================
\section{A Research Agenda for LLMs in Repeated Public Goods Games}
\label{sec:agenda}
% =====================================================================

The gaps above suggest a coherent programme. We organise it as a design space
plus a set of concrete, novel directions.

\paragraph{A design space for repeated-PGG agent experiments.}
A study can be positioned along independent axes, most of which are currently
under-sampled: (i) \emph{horizon} (few rounds vs.\ hundreds; known vs.\ unknown
endpoint, to expose endgame effects); (ii) \emph{enforcement} (none / peer
punishment / reward / pool punishment / institutional choice
\citep{gurerk2006competitive,sigmund2010social}); (iii) \emph{information}
(anonymous vs.\ persistent identities and reputation); (iv) \emph{communication}
(none vs.\ free-form negotiation, shown to be pivotal in \citep{piatti2024govsim});
(v) \emph{population} (single model, mixed models, LLM--human mixes); (vi)
\emph{payoff structure} (linear vs.\ threshold/collective-risk vs.\ commons); and
(vii) \emph{memory} (context window, summarisation, external memory).

\paragraph{Promising and largely novel directions.}
\begin{enumerate}
  \item \textbf{Long-horizon dynamics and benchmarking against human stylized
  facts.} Run PGGs over many rounds and test directly for decay, conditional
  cooperation, and endgame effects against Table~\ref{tab:facts}; only isolated
  studies yet reach this regime \citep{liu2026memory}. A standardised, open
  repeated-PGG benchmark (fixed protocol, metrics, human baselines)---in the spirit
  of emerging mechanism benchmarks \citep{tewolde2026coopeval}---would itself be a
  substantial contribution.
  \item \textbf{A strategy-method characterisation of LLM conditional
  cooperation.} Adapt \citet{fischbacher2001conditional} to elicit each model's
  contribution as a function of others' average contribution, classifying models
  as conditional cooperators, free-riders, or other types.
  \item \textbf{Mechanistic study of the reasoning--cooperation trade-off.}
  Causally isolate why reasoning models defect \citep{guzman2025corrupted} by
  varying chain-of-thought, decoding, persona, and alignment, and test
  interventions that restore cooperation.
  \item \textbf{Reputation and indirect reciprocity with persistent identities.}
  Building on first gossip-based studies \citep{cross2025validating,zhu2026gossip},
  re-introduce persistent identity and reputation into the institutional PGG and
  test whether LLM cooperation responds as human cooperation does
  \citep{milinski2002reputation,rockenbach2006efficient}.
  \item \textbf{The reward--punishment asymmetry and second-order free-riding.}
  Systematically probe why LLMs avoid punishment, whether they sustain enforcement,
  and whether metanorms \citep{horiguchi2024evolution} stabilise it.
  \item \textbf{Threshold and collective-risk PGGs with LLMs.} Port the
  climate-style dilemma \citep{milinski2008collective,santos2011risk} to LLM
  agents and measure risk sensitivity, threshold behaviour, and the effect of
  group size. Existing work either uses a single model or has models author
  strategies \citep{slumbers2025port,willis2026hundreds}; a systematic, multi-model
  study of agents \emph{deliberating} round by round would occupy a still-open,
  high-impact niche.
  \item \textbf{Evolutionary analysis of LLM PGG strategies.} Where
  \citet{willis2025cooperate} and \citet{willis2026hundreds} have models
  \emph{author} strategies, complement this by extracting the effective strategies
  of agents that \emph{play} the $N$-person PGG turn by turn, then study their
  evolutionary stability and invasibility under replicator dynamics.
  \item \textbf{Heterogeneous and mixed-population societies.} Compose groups of
  different model families, and of LLMs with humans, to study cooperation,
  exploitation, and convention formation \citep{ashery2025emergent} in realistic
  mixed settings.
  \item \textbf{Cooperation as an AI-safety property.} Frame multi-agent
  cooperation in repeated PGGs as an alignment target, connecting the empirical
  findings to the governance of deployed multi-agent systems
  \citep{guo2024multiagent,sun2025gametheory}.
\end{enumerate}

\noindent
Taken together, these directions share a common logic: the human and evolutionary
PGG literature provides a rich, quantitative set of hypotheses about \emph{repeated}
cooperation, while current LLM studies have tested only a small corner of it. The
clearest novelty therefore lies in (a) genuinely long-horizon experiments, (b)
designs that re-introduce the mechanisms humans rely on (punishment, reputation,
risk), and (c) bringing evolutionary and benchmarking rigour to LLM-authored
strategies.

% =====================================================================
\section{Conclusion}
\label{sec:conclusion}
% =====================================================================

The \pgg{} has been, for half a century, the sharpest lens on the conflict
between private incentive and collective good, and its repeated form is where the
mechanisms that rescue cooperation---reciprocity, reputation, punishment, norms,
and institutions---actually operate. As large language models become interacting
agents, the same lens is now being turned on them. The emerging evidence is that
capable LLMs cooperate substantially and reproduce many human regularities, that
their behaviour is highly model- and prompt-dependent, that reasoning optimisation
can paradoxically erode cooperation, and that they often achieve cooperative
outcomes through non-human mechanisms (reward over punishment, unforgiving
reciprocity). Yet the studies establishing these points have so far explored only
a narrow slice of the classical design space, and---crucially for anyone interested
in \emph{repeated} cooperation---the long-horizon dynamics that define the
repeated PGG remain largely uncharted. That gap, mapped in detail above, is also
the opportunity: a disciplined programme that takes the repeated PGG seriously, at
realistic horizons, with the mechanisms humans rely on and the rigour of
evolutionary game theory, is well positioned to produce genuinely novel results on
how---and whether---populations of language models can sustain the common good.

% =====================================================================
\bibliography{references}
% =====================================================================

\end{document}
